En una experiència, enviem radiació ultraviolada contra una placa de plom i produïm efecte fotoelèctric. Els electrons que es desprenen de la placa són frenats totalment per una diferència de potencial elèctric que depèn de la longitud d'ona de la radiació ultraviolada incident. A partir de les mesures efectuades sabem que quan la longitud d'ona és $1,50\times 10^{-7}\ \mathrm{m}$, la diferència de potencial que frena els electrons és de 4,01~V, i quan la longitud d'ona és $1,00\times 10^{-7}\ \mathrm{m}$, la diferència de potencial de frenada és de 8,15~V. Calculeu:

\begin{apartat}{1}
Per a cada longitud d'ona, la velocitat màxima amb què els electrons són extrets de la placa de plom.
\end{apartat}

\begin{apartat}{1}
L'energia mínima (funció de treball) necessària per a extreure un electró de la placa de plom. Determineu la constant de Planck a partir d'aquestes dades.
\end{apartat}

\dades{%
  $Q_\text{electró} = -1,60\times 10^{-19}\ \mathrm{C}$\\
  $c = 3,00\times 10^{8}\ \mathrm{m\,s^{-1}}$\\
  $m_\text{electró} = 9,11\times 10^{-31}\ \mathrm{kg}$}
